\cotefile{60}
\scanpage{36}
\begin{center}6 - 31 -\hfill n\textsuperscript{o} 392\end{center}

que $\operatorname{div}_B(Z)$ n'est pas l'image canonique d'un diviseur de $D(A)$. (Raisonner par l'absurde en supposant que $e(\underline{P}/\underline{p})=1$ pour $\underline{P}=BZ$, $\underline{p}=P\cap A$ ; il y aurait alors dans $A$ une uniformisante de $B_{\unsure{P|\underline{P}}}$, nécessairement de la forme $a(X,Y,Z)Z$ avec $b=a(X,Y,0)\ne 0$ ; en déduire que l'on aurait $\Delta b/b=-XY$ et obtenir une contradiction en calculant $\Delta(-XY)$).

24)\quad a) Soient $L$ un corps de caractéristique $2$, $\Delta$ une dérivation de $L$, $K$ le sous-corps de $L$ noyau de $\Delta$ ; on suppose que $[L:K]=2$. Montrer que l'on a $\Delta^2=a\Delta$ avec $a\in K$, et que pour qu'un élément $t\in L$ soit de la forme $\Delta x/x$, il faut et il suffit que $\Delta t=at+t^2$.

b) Soient $B$ un anneau local factoriel de caractéristique $2$, $\underline{m}$ son idéal maximal, $L$ son corps des fractions, $\Delta$ une dérivation de $L$ ; on suppose que le sous-corps $K$ de $L$, noyau de $\Delta$, est tel que $[L:K]=2$ ; en outre, on suppose qu'il existe deux éléments $x$, $y$ de $\underline{m}$ tels que $\Delta x$ et $\Delta y$ engendrent l'idéal $\underline{q}$ de $B$ engendré par $\Delta(B)$. Montrer alors que si $t=\Delta z/z$ appartient à $\underline{q}$, il existe un élément inversible $u$ de $B$ tel que $t=\Delta u/u$, (écrire $t=r\Delta x+s\Delta y$ avec $r$, $s$ dans $B$ et utiliser a)).

25) Soient $k$ un corps de caractéristique $2$, $B$ l'anneau de séries formelles $k[[X,Y]]$, $L$ son corps des fractions, $\Delta$ la $k$-dérivation de $L$ définie par $\Delta(X)=Y^{2j}$, $\Delta(Y)=X^{2i}$, ($i$, $j$ entiers $\geq 0$) ; le sous-anneau $A$ de $B$ formé des $x\in B$ tels que $\Delta x=0$ est l'anneau des séries formelles de $k[[X,Y,Z]]$ où on substitue $X^2$ à $X$, $Y^2$ à $Y$ et $X^{2i+1}+Y^{2j+1}$ à $Z$.

a) Montrer que le groupe $C(A)$ contient un espace vectoriel sur $k$ de dimension $N(i,j)$, égale au nombre des couples d'entiers $(a,b)$ tels que $0\leq a<i$, $0\leq b<j$ et $(2j+1)a+(2i+1)b\geq 2ij$. (Utilisant l'exerc. 23 b) et l'exerc. 24 a), noter que les éléments de $L$ sont les séries formelles $F\in B$ telles que $\Delta F=F^2$ ; en attribuant à $X$ le poids $2j+1$ et à $Y$ le poids $2i+1$, décomposer $F$ en somme infinie de polynômes isobares ; si $L_q$ est le sous-groupe de $L$ formé des $F\in L$ dont les composantes isobares sont de poids $\geq q$, $L/L'$ est isomorphe à la somme directe des groupes $C_q/C_{q+1}$, où $C_q=L_q/(L'\cap L_q)$ ; calculer ces groupes pour $q\geq 4ij$, en utilisant l'exerc. 24 b)).

b) Montrer que l'idéal $AX^2$ de $A$ est premier ; si $A'=A[X^{-2}]$, en déduire que $C(A')$ et $C(A)$ sont isomorphes, (n\textsuperscript{o}\quad, prop.\quad). Montrer que $A'$ est un anneau de Dedekind, (considérer l'anneau $B'=B[X^{-2}]$, qui est entier sur $A'$ et est principal, et utiliser le chap. V, \S\ 2, n\textsuperscript{o} 4, th. 3). En déduire un exemple d'anneau de Dedekind dont le groupe des classes d'idéaux est infini.

\margin{35}

